\section{Order and the Soul}

Picking up on the theme of Order and Chaos from the previous post, I thought it might be of interest to explicate Plato's concept of an ordered soul in \emph{The Republic}.

For Plato, the soul (\emph{psyche}) is composed of three main faculties -- the \emph{epithumetikon}, the \emph{thumoeides}, and the \emph{logistikon}.

The \emph{epithumetikon} is literally the `little desirer' and is ``concerned with eating, drinking, sex, and all their followers.'' This is often translated as the appetitive faculty.

The \emph{thumoeides} is a higher faculty than the \emph{epithumetikon}, as it is ``always wholly set on mastery, victory, and good reputation.'' Whereas the `little desirer' wants material things and physical pleasures, the \emph{thumoeides} (sometimes translated as the `spirited' faculty though this is potentially misleading given that Spirit has different meanings in different contexts) desires immaterial qualities.

Both \emph{epithumetikon} and \emph{thumomeides} are derived from \emph{thumos}, which in Homeric times referred to the motive force in man, and which is traditionally associated with the chest (hence the naming of the thymus gland.) Plato says that these faculties correspond to ``money and opinion.''\footnote{If we were to transpose these terms into yogic terminology, we might look at the similarities between the tripartite soul and the three \emph{gunas}, and locate the \emph{epithumetikon} in the region of the genitals and stomach, and the \emph{thumoeides} in the chest, considered the seat of courage as in one who `has great heart.'}

The third and highest faculty is the \emph{logistikon}, which derives from \emph{logos}. This faculty desires learning and wisdom, which ultimately lead to the Absolute. Though Plato does not say so explicitly, we can reasonably infer that the \emph{logistikon} comprises the two higher types of consciousness of which he writes later on. The first is \emph{dianoia}, which is discursive reasoning and is related to mathematics. The second and higher faculty is \emph{nous}, which is direct perception of true reality. These later became \emph{ratio} and \emph{intellectus} in the philosophy of the Middle Ages.

Different types of men are ruled by different parts of the soul. A man governed by the appetitive faculty will be primarily concerned with mundane affairs and satisfying as many of his desires as he is able. There are many men like this.

A man governed by the \emph{thumoeides} will be more concerned with honor and mastery and fame, whether in his own eyes or merely in the eyes of others. Depending on what kind of honor he desires, such a man might be a wealthy businessman, or a famous actor, or a courageous soldier. These men are rarer.

According to Plato, the philosopher must be ruled by the \emph{logistikon} and its desire for learning and wisdom, and must properly order his soul so that the two lower faculties of the soul are ruled by this desire, rather than obscuring or obstructing it. This proper ordering of the soul is called \emph{sophrosune}, translated as self-control or moderation. Plato calls it a kind of harmony between all the parts of the soul, which ``\emph{becomes entirely one from the many}, moderate and harmonized. Then, and only then, he acts.''

\begin{quotex}
Therefore, when all the soul follows the philosophic and is not factious, the result is that each part may, so far as other things are concerned, mind its own business and be just and, in particular, enjoy its own pleasures, the best pleasures, and, to the greatest possible extent, the truest pleasures.
\end{quotex}


Note that Plato does not call for the suppression of the lower faculties. The proper order of the soul is not a tyranny, but a monarchy.

This proper ordering of the soul is the precondition for the practice of philosophy; that is, awakening the \emph{nous} and attaining true \emph{gnosis}. And this proper ordering cannot be achieved without taking up the ancient maxim \emph{Know Thyself}.


\flrightit{Posted on 2011-09-21 by Will}

\begin{center}* * *\end{center}

\begin{footnotesize}
\begin{sffamily}
\texttt{Graham on 2011-09-21 at 16:11 said:}

It's good to see you back. I find it strange that the governing virtue sophrosune (harmony of soul) is translated as moderation or self-control, concepts which seem rather restricted. Isn't justice the virtue of man as such, and moderation the virtue special to the lowest part of the soul?

Evagrius Ponticus said that the conscious state of the just man was apatheia, or detachment. This state he held to be prerequisite for contemplation. I appreciate what you're doing with Plato, Will, since it's essential to reconstructing a Western esotericism.

\hfill

\texttt{Will on 2011-09-23 at 06:48 said:}

Graham, Thank you for your kind words. I agree that `moderation' is a bad translation of sophrosune, and `self-control' is not much better. Like arete -- `excellence' or `virtue' -- it has no exact English equivalent.

I'm not sure how Plato's system or terminology fits in with later Christian writers such as the one you mention, who I am not familiar with. However, in regard to justice as the virtue of man as such, Plato would say that only the philosopher -- by which he means something more akin to jnani than any modern sense of the word `philosopher'- is truly just, because only he perceives reality as it is.

One could make a case for Plato's agreement with the notion that apatheia is a prerequisite for contemplation, or the practice of philosophy, insofar as he writes of the need to govern the passions of the lower parts of the soul. But, on the other hand, he writes of the philosopher as having a deep love of learning and an insatiable appetite for wisdom and truth, which suggests that he is not entirely without pathos, but rather guided by one dominant drive.

I think a key aspect of sophrosune is that the soul becomes unified and capable of acting as a unity, i.e. the philosopher is not pulled to and fro by conflicting desires and drives. Gurdjieff also stressed the importance of achieving this state. But sophrosune is a means, not an end. For Plato, the aspiring philosopher harmonizes his soul so as to be able to practice philosophy -- to search for wisdom and truth with the entirety of his being.

This is consistent with Traditional teachings, such as the Buddhist story of the man who asks a teacher what he has to do to attain enlightenment. The teacher holds the man under water until he is about to drown, then finally lets him up for air and tells him that when he wants enlightenment as badly as he wanted to breathe just now, he'll attain it.


\hfill

\end{sffamily}
\end{footnotesize}

